\documentclass[../main.tex]{subfiles}
\ifSubfilesClassLoaded{\setcounter{chapter}{13}}{}
\begin{document}

\chapter{Struktur Narasi di Notebook; Ringkasan untuk Pemangku Kepentingan}

\begin{subcpmk}
  \item \textbf{Sub-CPMK-P10:} Menyusun ringkasan eksekutif (teks + visual) agar pembaca non-teknis memahami temuan (RPS Mg.\ 13).
\end{subcpmk}

\SesiRPS{13}{P10}{$\sim$170 menit}{CPMK-P4}

\noindent CPMK-P4 menuntut komunikasi \textit{insight}: pembaca manajerial tidak wajib memahami setiap baris kode, tetapi harus memahami apa yang data dukung, seberapa yakin kita, dan apa implikasi tindakannya.

\section{Kemampuan akhir sesi}

Menyusun struktur Markdown berlapis dari ringkasan eksekutif hingga rekomendasi; menyematkan minimal satu visual dan satu tabel metrik yang mendukung argumen utama.

Gunakan bahasa yang menghindari jargon berlebihan; jika istilah teknis perlu disebut, beri analogi satu kalimat atau rujuk ke bagian lampiran teknis di notebook.

\section{Teori}

\begin{konsep}
\textit{Data storytelling} menghubungkan bukti (tabel, plot, metrik) dengan rekomendasi. Struktur umum: konteks bisnis $\rightarrow$ pertanyaan $\rightarrow$ metode ringkas $\rightarrow$ temuan utama $\rightarrow$ keterbatasan $\rightarrow$ tindakan yang disarankan.
\end{konsep}

Ringkasan eksekutif sebaiknya dapat dibaca mandiri dalam satu layar: pembaca memutuskan apakah akan menggulir ke bagian metode dan lampiran bukti. Plot dipilih yang mendukung satu klaim utama, bukan menampilkan semua eksperimen sekaligus.

\begin{figure}[htbp]
\centering
\begin{tikzpicture}[node distance=8mm]
  \node[dsboxwide, fill=cyan!8] (a) {Bukti\\plot/metrik};
  \node[dsboxwide, fill=cyan!8, right=of a] (b) {Temuan\\kalimat inti};
  \node[dsboxwide, fill=cyan!8, right=of b] (c) {Aksi\\rekomendasi};
  \draw[dsarrow] (a) -- (b);
  \draw[dsarrow] (b) -- (c);
\end{tikzpicture}
\caption{Rangka narasi ringkasan eksekutif: bukti kuantitatif mendukung temuan, temuan mengarah pada rekomendasi yang dapat ditindaklanjuti.}
\end{figure}


\begin{table}[htbp]
\centering
\caption{Struktur bagian narasi di notebook (disarankan).}
\label{tab:struktur-narasi}
\begin{tabular}{@{}p{4.2cm}p{8.5cm}@{}}
\toprule
\textbf{Bagian} & \textbf{Isi} \\
\midrule
Ringkasan eksekutif & Masalah, jawaban singkat, rekomendasi. \\
Metode (ringkas) & Sumber data, model utama, validasi. \\
Temuan & Bullet + rujuk gambar/tabel. \\
Risiko & Kualitas data, asumsi, etika. \\
\bottomrule
\end{tabular}
\end{table}

\section{Contoh kerangka Markdown}

Gunakan sel Markdown berurutan (bukan sebagai listing Python). Kerangka berikut hanya contoh judul seksi; isi tiap seksi dengan paragraf dan bukti (angka, plot) dari notebook Anda.

\noindent\textbf{Mengapa blok ini?} Heading Markdown (\texttt{\#\#}) memecah ``cerita'' analitik menjadi blok yang bisa diekspor ke slide atau PDF; urutan mengikuti alur CRISP-DM ringkas.

\begin{verbatim}
## Ringkasan eksekutif (5-7 kalimat)
## Pertanyaan analitik
## Data dan kualitas
## Model dan validasi (tanpa rumus panjang)
## Temuan utama (3 bullet + 1 gambar)
## Risiko dan keterbatasan
## Rekomendasi
\end{verbatim}

\section{Contoh kode (plot untuk slide ringkasan)}

Visual untuk manajemen cenderung sederhana dan berlabel jelas; hindari palet yang sulit dibaca saat diproyeksikan.

\noindent\textbf{Mengapa blok ini?} Diagram batang membandingkan sedikit kategori (model) dengan satu metrik (F1); angka di sini hanya ilustrasi---ganti dengan metrik hasil validasi Anda.

\begin{lstlisting}[caption=Plot sederhana untuk disematkan di narasi]
import matplotlib.pyplot as plt
vals = [0.82, 0.88, 0.91]  # skor contoh; ganti dari eksperimen nyata
labs = ["Model A", "Model B", "Model C"]  # label sumbu-x
plt.bar(labs, vals)  # diagram batang vertikal
plt.ylabel("F1 (contoh)")
plt.title("Perbandingan model (ilustrasi)")
plt.show()
\end{lstlisting}

\noindent\textbf{Mengapa blok ini?} \texttt{savefig} menulis figur aktif ke disk; \texttt{dpi} mengontrol resolusi cetak/proyektor; \texttt{bbox\_inches="tight"} memangkas margin putih berlebih.

\begin{lstlisting}[caption=Export figur ke berkas]
plt.savefig("fig_ringkasan.png", dpi=150, bbox_inches="tight")  # jalankan setelah plot dibuat di sel sebelumnya
\end{lstlisting}

\noindent\textbf{Mengapa blok ini?} \texttt{to\_markdown} menghasilkan tabel Markdown dari DataFrame---cocok untuk menempelkan ringkasan metrik ke README atau laporan tanpa mengetik tabel manual.

\begin{lstlisting}[caption=Tabel ringkasan metrik ke Markdown via print]
import pandas as pd
dfm = pd.DataFrame({"model": ["A", "B"], "f1": [0.8, 0.86]})  # contoh dua baris
print(dfm.to_markdown(index=False))  # index=False: sembunyikan indeks 0,1 di output
\end{lstlisting}

\section{Referensi}

\begin{itemize}[leftmargin=*]
  \item IBM CRISP-DM (komunikasi hasil). \url{https://www.ibm.com/docs/en/spss-modeler/saas?topic=dm-crisp-help-overview}
\end{itemize}

\begin{asesmen}
\textbf{Formatif:} satu notebook ``cerita'' dari hasil minggu 7--12 dengan struktur di atas.
\end{asesmen}

\begin{rangkuman}
Minggu 13 mengalihkan fokus dari kode ke komunikasi keputusan berbasis data.

\textbf{Kata kunci:} ringkasan eksekutif, visual pendukung
\end{rangkuman}

\end{document}
